\documentclass[12pt]{article}

% -------------------------------------------------
% PAGE & LANGUAGE
% -------------------------------------------------

\usepackage[a4paper,margin=2cm]{geometry}

% -------------------------------------------------
% MATHEMATICS
% -------------------------------------------------

\usepackage{amsmath}
\usepackage{amssymb}
\usepackage{amsthm}
\usepackage{mathtools}
\usepackage{yhmath}

% Physics notation
\usepackage{physics}

% -------------------------------------------------
% GRAPHICS & PLOTS
% -------------------------------------------------

\usepackage{graphicx}
\usepackage{tikz}
\usepackage{pgfplots}

\pgfplotsset{compat=1.18}
\graphicspath{{./images/}}

% -------------------------------------------------
% COLORS & BOXES
% -------------------------------------------------

\usepackage[dvipsnames,svgnames,x11names]{xcolor}
\usepackage{tcolorbox}
\tcbuselibrary{skins,breakable}
\newtcolorbox{demobox}[1][]{
    enhanced,
    breakable,
    colback=black!2,
    colframe=black!70,
    boxrule=0pt,
    borderline west={2pt}{0pt}{black!70},
    sharp corners,
    left=10pt,
    right=7pt,
    top=6pt,
    bottom=6pt,
    before skip=12pt,
    after skip=12pt,
    fonttitle=\bfseries,
    title={Demostració},
    #1,
    before upper={\setlength{\parskip}{7pt}}

}

\newtcolorbox{examplebox}[1]{
    enhanced,
    breakable,
    colback=black!1,
    colframe=black!15,
    boxrule=0.5pt,
    arc=2mm,
    borderline west={2pt}{0pt}{black!55},
    left=12pt,
    right=12pt,
    top=10pt,
    bottom=10pt,
    before skip=12pt,
    after skip=12pt,
    title={\textbf{Exemple #1}},
    fonttitle=\normalsize,
    coltitle=black,
    attach boxed title to top left={
        xshift=8pt,
        yshift=-\tcboxedtitleheight/2
    },
    boxed title style={
        colback=white,
        colframe=white,
        boxrule=0pt,
        left=3pt,
        right=3pt,
        top=1pt,
        bottom=1pt
    },
    before upper={\setlength{\parskip}{7pt}}
}

% -------------------------------------------------
% DOCUMENT FORMATTING
% -------------------------------------------------

\usepackage{enumitem}
\usepackage[expansion=false]{microtype}
\usepackage{booktabs}
\usepackage{float}
\usepackage{ragged2e}

% -------------------------------------------------
% CHEMISTRY
% -------------------------------------------------

\usepackage[version=4]{mhchem}

% -------------------------------------------------
% LINKS
% --------------------------------------------------

\usepackage[colorlinks=true,linkcolor=Black,urlcolor=Black,citecolor=Black]{hyperref}
\usepackage{tocloft}
\usepackage{cite}

% -------------------------------------------------
% CUSTOM COMMANDS
% -------------------------------------------------

\newcommand{\R}{\mathbb{R}}
\newcommand{\N}{\mathbb{N}}
\newcommand{\Q}{\mathbb{Q}}
\newcommand{\Z}{\mathbb{Z}}

% -------------------------------------------------
% DOCUMENT STYLE
% -------------------------------------------------

\linespread{1.25}

\setlength{\parskip}{0.8em}
\setlength{\parindent}{0pt}

\setlength{\heavyrulewidth}{0.4pt}
\setlength{\heavyrulewidth}{0.4pt}
\setlength{\lightrulewidth}{0.4pt}
\setlength{\cmidrulewidth}{0.4pt}

\title{\textbf{Apunts Àlgebra}}
\author{\textit{Gerard Cosp Lores}}
\date{\today}

% -------------------------------------------------
% TABLE OF CONTENTS
% -------------------------------------------------

\renewcommand{\contentsname}{Índex}

% Title
\renewcommand{\cfttoctitlefont}{%
    \Huge\bfseries\color{Black}
}

% Line underneath title
\renewcommand{\cftaftertoctitle}{%
    \par\vspace{0.3cm}
    \noindent\color{Black}\rule{\textwidth}{1.2pt}
    \vspace{0.4cm}
}

% Section appearance
\renewcommand{\cftsecfont}{%
    \large\bfseries\color{Black}
}

\renewcommand{\cftsecpagefont}{%
    \large\bfseries\color{Black}
}

% Subsection appearance
\renewcommand{\cftsubsecfont}{%
    \normalsize
}

\renewcommand{\cftsubsecpagefont}{%
    \normalsize\color{Black}
}

% Spacing
\setlength{\cftbeforesecskip}{10pt}
\setlength{\cftbeforesubsecskip}{3pt}

% Dotted leaders
\renewcommand{\cftsecleader}{%
    \color{Black}\cftdotfill{\cftdotsep}
}

\renewcommand{\cftsubsecleader}{%
    \color{Black}\cftdotfill{\cftdotsep}
}

% Number spacing
\setlength{\cftsecnumwidth}{2.5em}
\setlength{\cftsubsecnumwidth}{3.5em}

% -------------------------------------------------

\begin{document}
\justifying

\maketitle
\tableofcontents

\newpage

\section{Definicions Bàsiques d'Algebra}

D'aqui en endavant anomenem a qualsevol cos \textbf{K}, un cos és aquella llista d'elements que tenen dues operacions bàsiques $(K, +, \cdot)$ on (al contrari que els anells) pots dvidir $\forall k \in K - 0$

\begin{itemize}
    \item + $K \cross K \rightarrow K$\\
    $(k_1, k_2) \rightarrowtail k_1 + k_2$
    \item $\cdot$ $K \cross K \rightarrow K$\\
    $(k_1, k_2) \rightarrowtail k_1 \cdot k_2$ 
\end{itemize}

Aquestes operacions compleixen les seves \textbf{propietats habituals a $\Q, \R$}

\textit{K amb +}
\begin{itemize}
    \item \textbf{Associativa}; $k_1 +(k_2 + k_3) = (k_1 + k_2) + k_3$
    \item \textbf{Commutativa}; $k_1 + k_2 = k_2 + k_1$
    \item \textbf{Element neutre}; (0) $k_i + 0 = k_i$
    \item \textbf{Element oposat}; $\forall k, \exists \overline{k} \mid k + \overline{k} = 0$
\end{itemize}

\textit{K amb $\cdot$}
\begin{itemize}
    \item \textbf{Associativa}; $k_1 \cdot(k_2 \cdot k_3) = (k_1 \cdot k_2) \cdot k_3$
    \item \textbf{Commutativa}; $k_1 \cdot k_2 = k_2 \cdot k_1$
    \item \textbf{Element neutre}; (1) $k_i \cdot 1 = k_i$
    \item \textbf{Element oposat}; $\forall k, \exists k^-1 \mid k \cdot k^-1 = 1$
\end{itemize}

Com a cossos tenim per exemple $f(x), \R, \Q, F_2 = \Z/2\Z$ on en l'ultim només existeix 0, 1 on $1 + 1 = 0$.

Anomenem \textbf{grup abelià} aquell cos que permet una operació que commuta com $(K, +), (K, \cdot)\dots$. Un conjunt de dos operacions és un \textbf{anell} per exemple $\R = (K, +, \cdot)$

\section{Matrius}

Una matriu és un cos K amb entrades $K^{m\cross n}$ que ens permet gestionar un array molt gran de nombres, soldrem usar $\Q, \R, \dots$

Dins de cada matriu es troba un element A = ($a_{ij}$) on $ 1 \leq i \leq m, 1 \leq j \leq n$. I hi ha diverses operacions que podem fer amb les matrius. Podem fer diverses operacions a les matrius de manera que les anomenem \textbf{transformacions fonamentals} ($\rightsquigarrow$)

\begin{itemize}
    \item Canviar una fila per una altra, $F_i \leftrightarrow F_j$
    \item Multipliquem una fila per un escalar $\lambda \in K$, $kF_i$
    \item Li sumem a una fila una altra multiplicada per un escalar, $F_i \rightarrowtail kF_j$
\end{itemize}

Qualsevol cosa es pot fer amb \textbf{columnes} de la mateixa manera. 

Una cosa important de les matrius i les transoformacions és que poden ser descrites com \textbf{matrius elementals} $E_i$ que són les transformacions lineals en una matriu identitat de manera que si fem $E_i \cdot A = A'$, obtenim la matriu A amb la transformació anterior. Si fem $A\cdot E_1 = A'$, hem dut a terme la transformació per les colmnes de la matriu.
\textbf{\textcolor{red}{Recordem una matriu $A\in K^{n\cross m}$ només pot ser multiplicada per una amb m files o n columnes al altre bando}.} De manera que si volem fer una transformació per la dreta o l'esquerra haurem de complir les condicions.

Per exemple;

\begin{itemize}
    \item \textbf{Canviem} $F_1 \leftrightarrow F_2$ = $ \begin{pmatrix}
0 & 1 & 0 \\
1 & 0 & 0 \\
0 & 0 & 1 
\end{pmatrix}  $

    \item \textbf{Multipliquem} la $F_1$ per $-3$ = $ \begin{pmatrix}
-3 & 0 & 0 \\
0 & 1 & 0 \\
0 & 0 & 1 
\end{pmatrix}  $

    \item \textbf{Multipliquem} la $F_1$ per $2$ i li \textbf{sumem} a $F_2$ = $ \begin{pmatrix}
1 & 0 & 0 \\
2 & 1 & 0 \\
0 & 0 & 1 
\end{pmatrix}  $
\end{itemize}

Una operació interesant és que \textbf{la matriu inversa} és l'operació contrària de la matriu $E^{-1}_1 \cdot E_1 = E_1 \cdot E^{-1}_1 = Id_n$.

Podem usar aquesta propietat per dur a terme una matriu $P = \prod_{i = 0}^{n} E_i = E_{1} \dots E_n$ on quan apliquem la matriu P es dona una serie de \textbf{transformacions lineals} aquesta la podem usar per trobar una matriu que multiplicada per P ens dona la forma \textbf{reduïda} de la matriu original.

Cal remarcar que les matrius $K^{n\cross n} \supset GL_n(K) = \{A \in K{n\cross n} | \exists A^{-1}\}$ Aixó vol dir que totes les matrius quadrades sense relacions lineals entre files o columnes tenen una inversa $AA^{-1} = A^{-1}A = Id_n$ de manera que $Rang(A) = n$

Si tenim una matriu al costat d'una identitat i apliquem uns canvis per obtenir $A'$ tindrem a l'identitat una matriu $P$ que \textbf{ens donara tot el conjunt de transformacions lineals que hem d'aplicar a A per obtenir A'}. Que pot ser la forma reduïda o la forma normal de Gauss-Jordan

\[
\left(
\begin{array}{ccc|ccc}
1 & 0 & 0 & a_{11} & a_{12} & a_{13}\\
0 & 1 & 0 & a_{21} & a_{22} & a_{23}\\
0 & 0 & 1 & a_{31} & a_{32} & a_{33}
\end{array}
\right)
\quad\rightarrowtail \quad
\left(
\begin{array}{ccc|ccc}
p_{11} & p_{12} & p_{13} & a'_{11} & a'_{12} & a'_{13}\\
p_{21} & p_{22} & p_{23} & a'_{21} & a'_{22} & a'_{23}\\
p_{31} & p_{32} & p_{33} & a'_{31} & a'_{32} & a'_{33}
\end{array}
\right)
\]

\subsection{Sistemes d'Equacions}

Anomenem sistema d'equacions a un \textbf{conjunt d'equacions} amb les mateixes incògnites per així poder donar una solució.

Un sistema qualsevol es pot expressar de dos maneres, la manera clàssica i la manera \textbf{matricial}.

\[
\left\{
\begin{aligned}
a_{11}x_1+a_{12}x_2+a_{13}x_3 &= b_1\\
a_{21}x_1+a_{22}x_2+a_{23}x_3 &= b_2\\
a_{31}x_1+a_{32}x_2+a_{33}x_3 &= b_3
\end{aligned}
\right.
\quad\longrightarrow\quad
\left(
\begin{array}{ccc|c}
a_{11} & a_{12} & a_{13} & b_1\\
a_{21} & a_{22} & a_{23} & b_2\\
a_{31} & a_{32} & a_{33} & b_3
\end{array}
\right)
\]

Anomenem \textbf{matriu ampliada ($A^*$)} a la matriu amb els termes independents, per resoldre qualsevol sistema \textbf{esglaonem la matriu} i la transformem a la manera clàssica per aïllar les incògnites.

Quan anem esglaonant ens podem trobar en 3 situacions;
\begin{itemize}
    \item $Rang(A)\neq Rang(A^*)$; SI
    \item $Rang(A) = Rang(A^*) = \#\,\text{incògnites}$; SCD
    \item $Rang(A) = Rang(A^*) < \#\,\text{incògnites}$; SCI
\end{itemize}

Si tenim un sistema on la matriu de coeficients i de termes independents és;
\[
\left(
\begin{array}{ccc|c}
    x\\
    y\\
    z\\
    .\\
    .\\
    .\\
\end{array}
\right)
\]\[
\left(
\begin{array}{ccc|c}
    0\\
    0\\
    0\\
    0\\
    0\\
    0\\
\end{array}
\right)
\]

Ens trobem amb un \textbf{sistema homogeni}, i n'hi ha dos tipus.

\begin{itemize}
    \item SCD amb solució $(0,0,\dots)$. - \textit{Incompatible}
    \item SCI amb infinites solucions. - \textit{Compatible}
\end{itemize}

Per resoldre un sistema també podem fer ús del mètode de cràmer;

\begin{demobox}

Donat un sistema lineal

$$
AX=B,
$$

si

$$
\det(A)\neq 0,
$$

el sistema té una única solució, que es pot obtenir mitjançant el \textbf{mètode de Cramer}:

$$
\boxed{x_i=\frac{\det(A_i)}{\det(A)}}
$$

on $A_i$ és la matriu que s'obté substituint la \textit{columna $i$ de $A$} per la columna $B$.

Per a un sistema de tres incògnites:

$$
\boxed{
x=\frac{\det(A_x)}{\det(A)},
\qquad
y=\frac{\det(A_y)}{\det(A)},
\qquad
z=\frac{\det(A_z)}{\det(A)}
}
$$

\end{demobox}

Si se'ns proposa un sistema amb uns paràmetres ($a \in R$) hem d'estudiar per rangs en funció del paràmetre.

Estudiem el $det(A)$ i trobem per quins a $det(A) = 0$ i per cadascun busquem el \textbf{rang de la matriu ampliada i comparem}.

\section{Espais Vectorials}

Fixem E-espai vectorial \textbf{(E-ev)}

Idea: Cos on es pot dur a terme \textbf{suma} i \textbf{multiplicació per escalar} de manera $(a_1, b_1) + (a_2, b_2) = (a_1 + a_2, b_1 + b_2)$ i ànalogament $\lambda(a_1, b_1) = (\lambda a_1, \lambda b_1), \lambda\in\R$.

\begin{itemize}
    \item + $E \cross E \rightarrow E$\\
    $(\vec{e}_1, \vec{e}_2) \rightarrowtail \vec{e_1} + \vec{e_2}$ (complint, Associativa, Commutativa)\\
    $\exists \vec{0_e} \mid \vec{e} + \vec{0_e} = \vec{e}, \forall \vec{e} \in E$\\
    $\exists \vec{\overline{e}} \mid \vec{e} + \vec{\overline{e}} = \vec{0_e}, \forall \vec{e} \in E$
    \item $\cdot \lambda$ $K \cross E \rightarrow E$\\
    $(k, \vec{e}) \rightarrowtail k \cdot \vec{e}$
    \begin{itemize}
        \item $(a_kb)_E \cdot u = a_k\cdot(b_E u)$
        \item $(a+b)_E \cdot u = a_E \cdot u + b_E \cdot u$
        \item $a_E(u+v) = a_E\cdot u + a_E \cdot + v$
        \item $1_E \cdot u = u$
    \end{itemize}
\end{itemize}

On es satisà que $u,v\in E$ i $a,b \in K$

\subsection{Subespai Vectorial}

Definim E un K-espai vectorial ($\R^2 i \R^3$) i un F K-espai vectorial tal que $F\subseteq E$ on es dona $(+, \cdot)$ i el resultant sempre ha d'estar dins de E.

% figura cercle amb suma fora del subespai

Com es veu en la figura si prenem 2 vectors $\vec{u} + \vec{v} = \vec{t} \notin E$. \textbf{\textcolor{red}{Qualsevol espai o qualsevo subespai ha de tenir el vector $\vec{0_{\R^2}}$}}. Necessitem el $\vec{0_{\R^2}}$ ja que no el podem obtenir aquest vector i el necessitem per sumar i restar.

Un subespai vectorial pot ser la recta que passa per $\vec{0_{\R^2}}$ de manera que es formen infinits vectors multiplicats per infinits escalars (creant aquesta recta). 

Una combinació lineal d'una familia de vectors / $<u_1, \dots u_n>\in E$ és un vector $u$ de E que s'obté multiplicant cada vector per un escalar $\lambda_n$ i sumant-los fent així; $<u_1,\dots,u_n>_K$ Totes les combinacions lineals dels vectors.

\textbf{IMPORTANT}: F és K-Subespai Vectorial $\leftrightarrow$

\begin{itemize}
    \item 1. $\forall \vec{u}, \vec{v}\in F, \text{ es té } u_E+v \in E$
    \item 2. $\forall \vec{u}\in F i \lambda \in K, \text{ es té } \lambda\vec{u}\in E$
\end{itemize}

Donat F tal que $<u_1,\dots, u_n>$, diem que aquests vectors són \textbf{un sistema de generadors per a F}, la base son els minims generadors d'un E.V, és a dir sense \textit{combinacions lineals} ni \textit{$\vec{0_{\R^2}}$}

Si diem que F és un Sub espai vectorial de E fix i \textbf{finit}, totes les $<f_1,\dots,f_n>_K = F \in E$ de manera que les \textbf{combinacions lineals d'aquests vectors formen l'espai}. Els casos més trivials d'un subespai vectorial, poden ser;

\begin{itemize}
    \item \textbf{Punts} $\vec{0}_{E}$
    \item Qualsevol \textbf{recta} que \textbf{passa per l'origen}
    \item \textbf{Tot} $E$
\end{itemize}

Cal remarcar que per tot $E$ un vector ($\vec{u} \in F \leftrightarrow u \in <\vec{u_1}, \dots, \vec{u_n}>_K$) és a dir que un vector pertany a un subespai si és combinacio lineal (\textbf{suma i multiplicació per escalar} $\vec{u} = \lambda_1\vec{u_1} + \dots + \lambda_n\vec{u_n}$).

\begin{demobox}
    Si $\vec{u} \in <\vec{u_1}, \dots, \vec{u_n}>_K$ es verifica que podem trobar $\vec{u}$ en $<\vec{u}, \vec{u_1}, \dots, \vec{u_n}>_K$ si $<\vec{u}, \vec{u_1}, \dots, \vec{u_n}>_K = <\vec{u_1}, \dots, \vec{u_n}>_K$ ja que podem trobar;

    \[\vec{u} = 1\vec{u} + 0 \vec{u_n}\dots\]

    Com que $\vec{u} \in <\vec{u_1}, \dots, \vec{u_n}>_K$, clarament cada combinació lineal podem donar un valor $\lambda_0 = 0$ de manera que ens queda que $<\vec{u}, \vec{u_1}, \dots, \vec{u_n}>_K = <\vec{u_1}, \dots, \vec{u_n}>_K$ i així es verifica que si tenim un espai vectorial on hi pertany un vector qualsevol \textbf{combinació lineal amb el vector també hi pertany}.
\end{demobox}

\begin{examplebox}{1}
    \[F = <(1,1,1), (1,2,3), (2,3,4)>\]

    Com que el vector $(1,1,1) \in <(2,3,4), (1,2,3)>_K$ és combinacio lineal dels altres, el podem eliminar del nostre sistema de generadors. Pel Lema de que una matriu un es fan operacions lineals \textbf{no perd la seva informació} podem escriure,

    \[A =\begin{pmatrix}
Fila_{1A} \\
\hline .\\
. \\
\hline \\
Fila _{nA} 
\end{pmatrix} \rightsquigarrow (\text{Operacions Elementals}) \begin{pmatrix}
Fila_{1A'} \\
\hline .\\
. \\
\hline \\
Fila _{nA'} / (0,\dots,0) 
\end{pmatrix}\]

    I reduïr la matriu a forma gauss jordan o esgalonada i veure quines files són combinació lineal de les altres per \textbf{simplificar el sistema de generadors}.
\end{examplebox}

Bàsicament com hem vist al exemple si tenim una matriu $M \in K^{m\cross n}$ si $Rang(A) < m$ és tracta d'un sistema de generadors amb \textbf{combinacions lineals}, si $Rang(A) = m$, tots els vectors del sistema de generadors són \textbf{linealment independents}.

\textit{Uns vectors són linealment independentment} si i només si \textbf{cap es pot construïr com a combinació dels altres} i en cas que igualem els vectors al vector origen ($\mu_1\vec{u_1}+ \dots + \mu_n\vec{u_n} = 0$) és un \textbf{SCD} homogeni amb úniques solucions ($\mu_1, \dots, \mu_n = 0$).

\subsection{Sistemes de Generadors i Bases}

Parlem d'un sistema de generadors tots els vectors (independentment si són C.L) dels altres o no, \textbf{que ens ajuden a construïr tot l'espai}. En el sistema de generadors podem afegir o eliminar vectors combinació lineal dels altres.

Com que un sistema de generadors també esta format per C.L dels altres \textbf{n'hi ha infinits sistemes generadors}.

\begin{examplebox}{2}
    Donat un sistema de generadors d'un espai vectorial E $\{\sin x, \cos x, 1\}, f = (\R, \R)^2$, volem veure quines solucions fan cert el sistema que descriu el sistema de generadors (posem en columnes els elements i l'element $0_{E}$). Volem veure si hi ha alguna possible combinació lineal entre els 3 elements que sigui zero sense que els coeficients $\mu_n$ siguin, de manera que;

    \[\mu_1 \sin x + \mu_2 \cos x + \mu_3 \cdot 1 = 0 \]

    Donem valors a les x i resolem;

    \begin{itemize}
        \item $x = 0 \therefore \mu_2 + \mu_3 = 0$
        \item $x = \frac{\pi}{2} \therefore \mu_1 + \mu_3 = 0$
        \item $x = \pi \therefore -\mu_2 + \mu 3 = 0$
    \end{itemize}

    Podem arribar a la conclusió que $\mu_1, \mu_2, \mu_3 = 0$ per tant com que el resultat són zeros no podem construïr cap element amb combinació dels altres.
\end{examplebox}

Per tant en aquest exemple 2 com que són \textbf{linealment independents}, $\{\sin x, \cos x, 1\}$ formen un \textit{sistema de generadors}.

\begin{demobox}
    Si $<\vec{u_1}, \dots, \vec{u_n}>$ de E són \textbf{linealment independents} es té que podem afegir $\vec{u} \notin <\vec{u_1}, \dots, \vec{u_n}>_K$ al nostra sistema de generadors de manera que no sigui \textbf{combinació lineal dels altres}.

    Demostrem això del contrari, si \textit{$\vec{u} \in <\vec{u_1}, \dots, \vec{u_n}>$} podem escriure;

    \[\vec{u} = \lambda_1\vec{u_1} + \dots + \lambda_n \vec{u_n}\]

    \[\vec{u} - \lambda_1\vec{u_1} + \dots + \lambda_n \vec{u_n} = \vec{0_E}\]

    De manera que podem transformar això en una equació;

    $x_0 \vec{u} -x_1 \vec{u_1} - \dots- x_n \vec{u_n} = 0$ que és un Sistema Homogeni amb solució $(0,\dots,0)$ per tant són \textbf{linealment depenents} si pertany.

    Però com que hem dit que no pertany \textbf{són Linealment Independents.}
\end{demobox}

Com podem reduïr un sistema de generadors? Doncs primer mirem quins són combinació lineal dels altres i ho reduim amb paràmetres als mínims múltiples possibles

\begin{examplebox}
    Siguin $E = <(1,2,3), (2,3,4), (3,4,5)>$ un sistema de generadors, volem reduïr i simplificar al màxim, creem un sistema tal que;

    \[x_1 \begin{pmatrix}
        1\\
        2\\
        3
    \end{pmatrix} + x_2\begin{pmatrix}
        2\\
        3\\
        4
    \end{pmatrix} + x_3 \begin{pmatrix}
        3\\
        4\\
        5
    \end{pmatrix} = \begin{pmatrix}
        0\\
        0\\
        0
    \end{pmatrix}\]

    \textbf{Col·loquem els vectors en files i ho igualem a $\vec{0_E}$} i obtenim el sistema;

    \[
    \begin{pmatrix}
        1 & 2 & 3 & 0 \\
        2 & 3 & 4 & 0 \\
        3 & 4 & 5 & 0 
    \end{pmatrix}
    \]

    Si resolem ens queda el sistema reduït com;

    \[
    \begin{pmatrix}
        1 & 1 & 1 & 0 \\
        0 & 1 & 2 & 0 \\
        0 & 0 & 0 & 0 
    \end{pmatrix}
    \]

    Si parametritzem queden dos vectors de la forma $(\lambda, -2\lambda, \lambda), (0, -2\lambda, 2\lambda)$ i si treiem $\lambda$ com a factor comú, ens queda que podem rescriure el nostre \textbf{sistema de generadors com;}

    \[<(0,-1,1), (1,-2,1)>\]
\end{examplebox}

Una \textbf{base} d'un espai E és un sistema de generadors on tots són \textit{linealment independents}, i és la familia màxima de vectors independents (màxims vectors independents, $\#vectors = dimensions$).

\begin{itemize}
    \item \textbf{Canònica}: Base habitual ($\vec{i}, \vec{j},\vec{k}$)
    \item \textbf{Ortogonal}: Els vectors formen $\alpha = 90$ entre ells
    \item \textbf{Normal}; Els vectors tenen $\abs{\vec{u}} = 1$
    \item \textbf{Ortonormal}; Els vectors tenen $\abs{\vec{u}} = 1$ i formen $\alpha = 90$ entre ells.
\end{itemize}

Si una base compleix tots els seus requisits de manera que tenim $B = <\vec{u_1}, \vec{u_2}, \vec{u_3}, \dots, \vec{u_n}>$ podem descomposar qualsevol vector $\vec{e}$ com;

\[\vec{e} = \lambda_1 \vec{u_1} + \lambda_2 \vec{u_2} + \lambda_3 \vec{u_3} + \dots + \lambda_n \vec{u_n}\]
\end{document}